\documentclass[11pt]{article}
\usepackage{amsmath,amssymb}
\usepackage[margin=1in]{geometry}
\usepackage{hyperref}
\emergencystretch=3em

\title{Lipschitz dependence of the limit coefficient on phase and amplitude}
\author{Claude, with Daniel Murfet}
\date{2026-09-07}

\begin{document}
\maketitle
\sloppy

\begin{abstract}
Unit 210 (programme A2, step 2) is the limit-side companion of unit 209. The functional
$(\xi,\eta)\mapsto\mathrm{phaseCoeff}(\xi,\eta)=D\int\eta(\pi u)J_p(\xi(\pi u))w(u)\,du$ of
Headline~XIII is Lipschitz in sup norm on bounded input sets:
\[
|\mathrm{phaseCoeff}(\xi,\eta)-\mathrm{phaseCoeff}(\xi',\eta')|
\le D\bigl(\|\eta-\eta'\|_\infty J_p(M)+\beta A\|\xi-\xi'\|_\infty J_{p+1}(M)\bigr)\int w
\]
for $|\xi|,|\xi'|\le M$, $|\eta'|\le A$ on the closed cube (\texttt{phaseCoeff\_sub\_le}), where
$J_{p+1}(M)=\int_0^\infty s^pe^{-\beta s^2+\beta Ms}\,ds$. The moment-level inputs are monotonicity
$J_p(a)\le J_p(M)$ for $a\le M$ and $|J_p(a)-J_p(a')|\le\beta|a-a'|J_{p+1}(M)$, both integrated from
the kernel bounds of unit~209. With \texttt{normalised\_lipschitz\_eventually} this makes the whole
family $\{F_N\}_{N\ge N_0}\cup\{F\}$ equi-Lipschitz on bounded input sets, the hypothesis of the
uniform-convergence step. Zero \texttt{sorry}/\texttt{axiom}; 9087 jobs.
\end{abstract}

\section{The things formalised}
{\small
\begin{verbatim}
theorem phaseMoment_le_of_le (hβ) (hp : 0 < p) (haM : a ≤ M) : phaseMoment β p a ≤
      phaseMoment β p M
theorem phaseMoment_sub_le (hβ) (hp) (ha : |a| ≤ M) (ha' : |a'| ≤ M) :
    |phaseMoment β p a - phaseMoment β p a'| ≤ β * |a - a'| * phaseMoment β (p + 1)
      M
theorem phaseCoeff_sub_le (hk) (hl) (hβ) (hmin) (continuity) (|ξ|,|ξ'| ≤ M) (|η'| ≤
      A)
    (|ξ - ξ'| ≤ δξ) (|η - η'| ≤ δη on closedCube d) :
    |phaseCoeff h k l β ξ η - phaseCoeff h k l β ξ' η'|
      ≤ 2 ^ (multCount (ratioExp h k) l - 1) * 2 * faceNorm h k l
          * ((δη * phaseMoment β (2 l) M + β * A * δξ * phaseMoment β (2 l + 1) M)
              * ∫ u in unitBox d, residualWeight h k l u)
\end{verbatim}
}

\section{Mathematics}
Pointwise on the face, $\eta J(\xi)-\eta'J(\xi')=(\eta-\eta')J(\xi)+\eta'(J(\xi)-J(\xi'))$ with
$J(\xi)\le J(M)$ and $|J(\xi)-J(\xi')|\le\beta|\xi-\xi'|J_{p+1}(M)$; the residual weight is
nonnegative, so the bound integrates. $s^{p-1}\cdot s=s^{(p+1)-1}$ turns the extra factor $s$ from
the kernel bound into the moment of order $p+1$ (\texttt{Real.rpow\_add\_one}).

\section{Lean notes}
\begin{itemize}
\item \texttt{gcongr} leaves the side goal $0\le J_{p+1}(M)$ when the atom is an opaque integral;
  discharge it with \texttt{all\_goals first | exact … | exact (phaseMoment\_pos …).le} rather than
  guessing the goal order.
\item \texttt{set v := faceProj h k l u} folds every occurrence and keeps the pointwise calculation
  readable; the bounds \texttt{hξM v hπ} apply because \texttt{faceProj\_mapsTo} lands in the
  closed cube.
\end{itemize}

\section{Next}
Unit 211: inputs as elements of $C([0,1]^d,\mathbb R)^2$ via the clamp extension
$u\mapsto(\max(0,\min(1,u_i)))_i$; equi-Lipschitz $\Rightarrow$ equicontinuous on bounded sets;
pointwise convergence $\Rightarrow$ uniform convergence on compact input sets. Unit 212: random
inputs, using tightness of the limit law in the Polish space $C([0,1]^d)^2$ and the small-probability
approximation lemma.
\end{document}
